% =========================================================================
% Frontmatter : résumé, table des matières, listes
% =========================================================================

\chapter*{Résumé}
\addcontentsline{toc}{chapter}{Résumé}

Ce mémoire présente la conception et l'implémentation d'un aquarium virtuel
animé écrit en C avec la bibliothèque graphique GTK4. Le système simule en
temps réel un écosystème de poissons multi-espèces à partir d'un modèle
comportemental dérivé des travaux de Reynolds sur le \emph{flocking}
(séparation, alignement, cohésion), enrichi de relations prédateur/proie,
d'un cycle de vie complet (naissance, mort, respawn) et d'une mécanique
d'alimentation basée sur la position de la \emph{bouche} des prédateurs.

L'apport principal du projet est la mise en place d'un \emph{système de
configuration externe} fondé sur le format TOML. Ce système permet à un
utilisateur non développeur d'ajouter une nouvelle espèce, d'en supprimer
une, ou d'ajuster finement n'importe quel paramètre de simulation
(vitesses, rayons de vol en banc, populations, prédateurs, etc.) en
modifiant un seul fichier texte~: \file{aquarium.conf}. Le système
gère les chemins de recherche standards (\file{./}, \file{\$XDG\_CONFIG\_HOME/},
\file{/etc/}), valide les références croisées entre espèces, fournit des
messages d'erreur lisibles, et autorise le rechargement à chaud
depuis l'interface graphique.

Sur le plan technique, le rapport détaille~: l'architecture modulaire
(modules \code{main}, \code{world}, \code{entity}, \code{species}, \code{conf},
\code{vec2})~; la boucle de simulation à pas fixe basée sur
\code{g\_timeout\_add}~; la gestion d'un pool d'entités à capacité bornée~;
la résolution des références d'espèces par index dynamique permettant
l'ajout d'espèces depuis le fichier de configuration sans recompilation~;
l'utilisation de transformations CSS GTK4 pour la rotation des sprites~;
et enfin l'intégration de la bibliothèque \code{tomlc99} comme parseur
TOML embarqué.

\textbf{Mots-clés~:} simulation comportementale, flocking, boids,
GTK4, langage C, configuration TOML, prédateur/proie, écosystème virtuel.

% --- Table des matières --------------------------------------------------
\tableofcontents

% --- Liste des figures ---------------------------------------------------
\listoffigures

% --- Liste des tableaux --------------------------------------------------
\listoftables

% --- Liste des codes (minted) --------------------------------------------
\listoflistings
